\documentclass{article}
\usepackage[T1]{fontenc}
\usepackage{lmodern}
\usepackage{graphicx}
\usepackage{amsmath,amssymb}
\usepackage[a4paper,margin=2cm]{geometry}
\usepackage[absolute,overlay]{textpos}
\setlength{\TPHorizModule}{1mm}
\setlength{\TPVertModule}{1mm}
\usepackage[indonesian]{babel}
\usepackage{hyperref}
\setlength{\emergencystretch}{2em}
% unit: Halaman 30

\begin{document}

% === Halaman 30 (python/native) ===

Dekripsi

• Dekripsi terhadap cipherteks merupakan kebalikan dari proses enkripsi. 

• DES menggunakan algoritma yang sama untuk proses enkripsi dan dekripsi. 

• Pada proses dekripsi urutan kunci yang digunakan adalah K16, K15, …, K1.

• Untuk tiap putaran 16, 15, …, 1, luaran pada setiap putaran deciphering adalah


Ri – 1 = Li 





             Li – 1 = Ri  f(Ri – 1, Ki) = Ri  f(Li, Ki)

30

\end{document}
